\documentclass[border=10pt]{standalone}
\usepackage{fontspec}
\setmainfont{Noto Serif}[Ligatures=TeXOff]
\XeTeXgenerateactualtext=1
\newfontfamily\bengalifont{Noto Serif Bengali}[Script=Bengali, NFSSFamily=bnfam]
\usepackage[Latin,Bengali]{ucharclasses}
% group-free font switching (as in header-common.tex): robust inside braces, cells and quotes
\makeatletter
\def\BN@saved{\rmdefault}
\def\BN@bn{bnfam}
\newcommand\BNin{\edef\BN@tmp{\f@family}\ifx\BN@tmp\BN@bn\else\let\BN@saved\BN@tmp\fi\fontfamily{bnfam}\selectfont}
\newcommand\BNout{\fontfamily{\BN@saved}\selectfont}
\makeatother
\setTransitionsFor{Bengali}{\BNin}{\BNout}
% Typographic quotes and dashes are not in a ucharclasses class (class 0), so a Bangla word right
% after “ or ‘ would not trigger the Bangla font. Treat them as Latin so that every switch fires.
\XeTeXcharclass"2018=\XeTeXcharclass`A
\XeTeXcharclass"2019=\XeTeXcharclass`A
\XeTeXcharclass"201C=\XeTeXcharclass`A
\XeTeXcharclass"201D=\XeTeXcharclass`A
\XeTeXcharclass"2013=\XeTeXcharclass`A
\XeTeXcharclass"2014=\XeTeXcharclass`A
\XeTeXcharclass"2026=\XeTeXcharclass`A
\usepackage{newunicodechar}
\newunicodechar{।}{{\bengalifont\char"0964}}
\newunicodechar{॥}{{\bengalifont\char"0965}}

\usepackage[table]{xcolor}
\usepackage{booktabs,tabularx,array}
\definecolor{bapilrule}{HTML}{1F2937}
\definecolor{bapilstripe}{HTML}{F3F4F6}
\renewcommand{\arraystretch}{1.28}
\setlength{\tabcolsep}{5pt}

\begin{document}
\begin{minipage}{17cm}
{\large\bfseries সারণি ৫২।}\ {\large এআই পণ্য ও সেবা ক্রয়ের আদর্শ শর্ত}\par\vspace{7pt}
\small
\rowcolors{2}{bapilstripe}{white}
\begin{tabularx}{\linewidth}{@{}>{\hsize=0.732\hsize\raggedright\arraybackslash}X>{\hsize=1.317\hsize\raggedright\arraybackslash}X>{\hsize=0.951\hsize\raggedright\arraybackslash}X@{}}
\arrayrulecolor{bapilrule}\toprule
\rowcolor{white}\textbf{শর্ত} & \textbf{প্রয়োজনীয়তা} & \textbf{উদ্দেশ্য} \\
\midrule
ডেটা-অবস্থান ও মালিকানা & সরকারি ডেটা বাংলাদেশে থাকবে ও সরকারের সম্পত্তি থাকবে; বিক্রেতার মডেল প্রশিক্ষণে ব্যবহার নয় & ডেটা ও সার্বভৌমত্ব সুরক্ষা (NFR-02) \\
মূল্যায়নের অধিকার & চুক্তির আগে ও পরে ক্রেতা BanglaPolicyQA ও রেড-টিম পরীক্ষায় পণ্য যাচাই করতে পারবেন & দাবির নয়, প্রমাণের ভিত্তিতে সিদ্ধান্ত \\
বাংলা কর্মদক্ষতার গ্রহণযোগ্যতা & গ্রহণের শর্ত সম্মত বাংলা ও ইংরেজি সীমা অর্জন & বাংলা প্রথম শ্রেণির ভাষা \\
নথিভুক্তি & প্রতিটি সংস্করণের সঙ্গে মডেল কার্ড, ডেটাশিট, SBOM ও জ্ঞাত সীমাবদ্ধতা প্রদান & স্বচ্ছতা ও নিরীক্ষা \\
স্থানান্তরযোগ্যতা ও প্রস্থান & উন্মুক্ত ফরম্যাট; সব ডেটা ও কনফিগারেশন রপ্তানি; চুক্তি শেষে রূপান্তর-সহায়তা & নির্ভরশীলতা এড়ানো (R10) \\
সোর্স কোড & সম্ভব হলে উন্মুক্ত সোর্স; অন্যথায় এসক্রো ও নিরীক্ষার অধিকার & দীর্ঘমেয়াদি রক্ষণাবেক্ষণযোগ্যতা \\
নিরাপত্তা ও ঘটনা & সারণি ৩২-এর নিয়ন্ত্রণ মেনে চলা; দুর্বলতা প্রকাশ; ২৪ ঘণ্টার মধ্যে ঘটনার বিজ্ঞপ্তি & নিরাপত্তা ও জবাবদিহি \\
পরিবর্তন নিয়ন্ত্রণ & মডেল বা বড় নির্ভরতা পরিবর্তনের আগে নোটিশ ও পুনর্মূল্যায়ন & স্থিতিশীল ও মূল্যায়িত আচরণ (G5) \\
প্রবেশযোগ্যতা & প্রমাণসহ WCAG 2.2 লেভেল AA সামঞ্জস্য & অন্তর্ভুক্তি (NFR-05) \\
সততা & স্বার্থের সংঘাত ঘোষণা; মূল্যায়ন বা শাসনে কোনো প্রভাব নয় & স্বাধীনতা \\
\bottomrule
\end{tabularx}
\par\vspace{3pt}{\scriptsize\color{gray} উৎস: BAPIL টেকনিক্যাল ব্লুপ্রিন্ট, সংস্করণ ১.০ (২০২৬)। CC BY 4.0।}
\end{minipage}
\end{document}
